% GENERATED FILE. Edit canonical lessons, then run npm run generate:books.
% canonical-source-hash: fnv1a64:b747f7520b317681
% canonical-lessons: ZH-C499-qiudui, ZH-C499-qiuyuan, ZH-C499-duizhang, ZH-C499-xunlian, ZH-C499-chuanqiu

\chapter{The Team}
\label{ch:zh-the-team}
\hlchaptermodality{499}

\begin{chapteropening}
\textbf{By the end of this chapter:} \emph{I can talk about a team, its players and captain, training and passing the ball.}

Complete the last lesson of chapter 499: I can talk about a team, its players and captain, training and passing the ball.
\end{chapteropening}

\section[qiúduì]{qiúduì — a sports team (in a ball game)}

\label{lesson:ZH-C499-qiudui}



\begin{quote}
Before the new one: say the Mandarin for to shoot at goal, then the Mandarin for to score a goal.
\end{quote}

\subsection*{You'll want to know: qiúduì}
\emph{qiúduì} — \textquotedblleft{}a sports team (in a ball game)\textquotedblright{}.

\begin{quote}
\emph{zhè shì wǒmen xuéxiào zuì hǎo de qiúduì} — This is our school's best team.
\end{quote}

Its characters come later; for now it is a word for the ear.

\begin{cousinweb}
\emph{qiúduì} is built from \zh{球} 'a ball' and \emph{duì} 'a troop', both inherited from Middle Chinese, \emph{gjuw} and \emph{dwojH} (Guangyun, in Baxter's transcription), which glosses the second 'a band of people'.
\end{cousinweb}

\begin{grammarlens}[title={What you've built}]
One more word for saying how it was.
\end{grammarlens}

\subsection*{Your turn}
\begin{itemize}
\raggedright
  \item \emph{Say it:} \emph{qiúduì}
  \item \emph{Say it:} \emph{qiúduì}, once more
  \item \emph{Say it:} \emph{zhè shì wǒmen xuéxiào zuì hǎo de qiúduì} — This is our school's best team.
  \item \emph{Recall:} say the Mandarin for to shoot at goal, then the Mandarin for to score a goal, then say \emph{qiúduì} again
\end{itemize}

\begin{tcolorbox}[breakable,colback=teal!4,colframe=teal!35!black,title={Before you move on}]
What does qiúduì mean? (\textquotedblleft{}A sports team (in a ball game)\textquotedblright{}.) Say it once more.
\end{tcolorbox}
\section[qiúyuán]{qiúyuán — a player in a ball game}

\label{lesson:ZH-C499-qiuyuan}



\begin{quote}
Before the new one: say the Mandarin for to score a goal, then the Mandarin for a sports team.
\end{quote}

\subsection*{You'll want to know: qiúyuán}
\emph{qiúyuán} — \textquotedblleft{}a player in a ball game\textquotedblright{}.

\begin{quote}
\emph{tā dāng le shí nián qiúyuán} — He was a player for ten years.
\end{quote}

Its characters come later; for now it is a word for the ear.

\begin{cousinweb}
\emph{qiúyuán} is built from \zh{球} 'a ball' and \emph{yuán} 'a member', both inherited from Middle Chinese, \emph{gjuw} and \emph{hjwen} (Guangyun, in Baxter's transcription), which glosses the second 'a count of things'; the sense 'a member' is not given there.
\end{cousinweb}

\begin{grammarlens}[title={What you've built}]
One more word for saying how it was.
\end{grammarlens}

\subsection*{Your turn}
\begin{itemize}
\raggedright
  \item \emph{Say it:} \emph{qiúyuán}
  \item \emph{Say it:} \emph{qiúyuán}, once more
  \item \emph{Say it:} \emph{tā dāng le shí nián qiúyuán} — He was a player for ten years.
  \item \emph{Recall:} say the Mandarin for to score a goal, then the Mandarin for a sports team, then say \emph{qiúyuán} again
\end{itemize}

\begin{tcolorbox}[breakable,colback=teal!4,colframe=teal!35!black,title={Before you move on}]
What does qiúyuán mean? (\textquotedblleft{}A player in a ball game\textquotedblright{}.) Say it once more.
\end{tcolorbox}
\section[duìzhǎng]{duìzhǎng — team captain; team leader}

\label{lesson:ZH-C499-duizhang}



\begin{quote}
Before the new one: say the Mandarin for a sports team, then the Mandarin for a player in a ball game.
\end{quote}

\subsection*{You'll want to know: duìzhǎng}
\emph{duìzhǎng} — \textquotedblleft{}team captain; team leader\textquotedblright{}.

\begin{quote}
\emph{wǒmen dōu xiǎng dāng duìzhǎng} — We all want to be captain.
\end{quote}

Its characters come later; for now it is a word for the ear.

\begin{cousinweb}
\emph{duìzhǎng} is built from \emph{duì} 'a troop' and \emph{zhǎng} 'a head', both inherited from Middle Chinese, \emph{dwojH} and \emph{trjangX} (Guangyun, in Baxter's transcription), which glosses the second 'great'; the sense 'a head of a group' is not given there.
\end{cousinweb}

\begin{grammarlens}[title={What you've built}]
One more word for saying how it was.
\end{grammarlens}

\subsection*{Your turn}
\begin{itemize}
\raggedright
  \item \emph{Say it:} \emph{duìzhǎng}
  \item \emph{Say it:} \emph{duìzhǎng}, once more
  \item \emph{Say it:} \emph{wǒmen dōu xiǎng dāng duìzhǎng} — We all want to be captain.
  \item \emph{Recall:} say the Mandarin for a sports team, then the Mandarin for a player in a ball game, then say \emph{duìzhǎng} again
\end{itemize}

\begin{tcolorbox}[breakable,colback=teal!4,colframe=teal!35!black,title={Before you move on}]
What does duìzhǎng mean? (\textquotedblleft{}Team captain; team leader\textquotedblright{}.) Say it once more.
\end{tcolorbox}
\section[xùnliàn]{xùnliàn — to do training (for a sport); to drill}

\label{lesson:ZH-C499-xunlian}



\begin{quote}
Before the new one: say the Mandarin for a player in a ball game, then the Mandarin for a team captain.
\end{quote}

\subsection*{You'll want to know: xùnliàn}
\emph{xùnliàn} — \textquotedblleft{}to do training (for a sport); to drill\textquotedblright{}.

\begin{quote}
\emph{wǒmen měi tiān xùnliàn liǎng ge xiǎoshí} — We train for two hours a day.
\end{quote}

Its characters come later; for now it is a word for the ear.

\begin{cousinweb}
\emph{xùnliàn} is built from \emph{xùn} 'to instruct' and \emph{liàn} 'to practise', both inherited from Middle Chinese, \emph{xjunH} and \emph{lenH} (Guangyun, in Baxter's transcription), which glosses the first 'to admonish' and the second 'white silk'; the sense 'to practise' is not given there.
\end{cousinweb}

\begin{grammarlens}[title={What you've built}]
One more word for saying how it was.
\end{grammarlens}

\subsection*{Your turn}
\begin{itemize}
\raggedright
  \item \emph{Say it:} \emph{xùnliàn}
  \item \emph{Say it:} \emph{xùnliàn}, once more
  \item \emph{Say it:} \emph{wǒmen měi tiān xùnliàn liǎng ge xiǎoshí} — We train for two hours a day.
  \item \emph{Recall:} say the Mandarin for a player in a ball game, then the Mandarin for a team captain, then say \emph{xùnliàn} again
\end{itemize}

\begin{tcolorbox}[breakable,colback=teal!4,colframe=teal!35!black,title={Before you move on}]
What does xùnliàn mean? (\textquotedblleft{}To do training (for a sport); to drill\textquotedblright{}.) Say it once more.
\end{tcolorbox}
\section[\texorpdfstring{\zh{传球} (chuánqiú)}{传球 (chuánqiú)}]{\zh{传球} (chuánqiú) — to pass the ball}

\label{lesson:ZH-C499-chuanqiu}



\begin{quote}
Before the new one: say the Mandarin for a team captain, then the Mandarin for to do training (for a sport).
\end{quote}

\subsection*{You'll want to know: \zh{传球}}
\textbf{\zh{传球}} — \emph{chuánqiú} — \textquotedblleft{}to pass the ball\textquotedblright{}.

\begin{quote}
\textbf{\zh{他传球很快}} \emph{tā chuánqiú hěn kuài} — He passes the ball quickly.
\end{quote}

\begin{cousinweb}
\textbf{\zh{传球}} (\emph{chuánqiú}) is built from \zh{传} 'to hand on' and \zh{球} 'a ball', both inherited from Middle Chinese, \emph{drjwen} and \emph{gjuw} (Guangyun, in Baxter's transcription), which glosses the first 'to turn'; the sense 'to pass on' it gives only under the other reading, \emph{drjwenH}.
\end{cousinweb}

\begin{grammarlens}[title={What you've built}]
One more word for saying how it was.
\end{grammarlens}

\subsection*{Your turn}
\begin{itemize}
\raggedright
  \item \emph{Say it:} \emph{chuánqiú}
  \item \emph{Say it:} \emph{chuánqiú}, once more
  \item \emph{Say it:} \emph{tā chuánqiú hěn kuài} — He passes the ball quickly.
  \item \emph{Recall:} say the Mandarin for a team captain, then the Mandarin for to do training (for a sport), then say \emph{chuánqiú} again
\end{itemize}

\begin{tcolorbox}[breakable,colback=teal!4,colframe=teal!35!black,title={Before you move on}]
What does \zh{传球} mean? (\textquotedblleft{}To pass the ball\textquotedblright{}.) Say it once more.
\end{tcolorbox}
